\begin{center}
\textbf{ABSTRACT}
\end{center}

$\!$\\

This report presents the dynamic modeling of a low-cost articulated robot, built with popsicle and skewer sticks and driven by a geared DC motor, that moves along a rope. Mass, dimensions and angles of each part were obtained with a scale, a caliper and photographs, and the inertia tensor of each component was computed from elementary homogeneous solids and composed with the parallel axis theorem, resulting in a 103.1~g robot whose center of mass lies 21.4~mm below the motor axis. The mechanism, in which two cranks 180$^\circ$ out of phase drive arms that act as levers about the hip axis and end in hooks, was modeled with the vector relations for velocity and acceleration of a rigid body and with the Newton-Euler equations projected onto the generalized speeds (Kane's method), checked against the Lagrange formulation with a relative difference of the order of $10^{-16}$. The contact between the hooks and the rope was treated as a pin, with hand exchange modeled as a perfectly plastic impact. The simulation with the motor at 30~rpm counterclockwise shows the robot moving forward, with a center of mass displacement of 131~mm in 8~s and an average speed of about 23~mm/s, against 43.5~mm/s predicted by pure kinematics, and with a vertical rope force that remains positive. The advance depends on the damping of the chassis swing, which was assumed, and on the adopted hook engagement rule.

\vspace{1cm}

\hspace{-1.3cm}Keywords: Rigid body dynamics, Newton-Euler, Inertia tensor, Articulated robot, Simulation.
